\chapter{Macular Hole}
\index{Macular Hole}
\label{sec:Macular Hole}

\section{Overview}
\begin{itemize}[noitemsep]
  \item A macular hole is a full-thickness defect in the center of the retina (the macula), which is responsible for sharp central vision
  \item It affects roughly 3 in 1000 people over 55 and is more common in women
  \item It causes progressive blurring and distortion of central vision and can lead to permanent vision loss if untreated
\end{itemize}
\section{Causes}
This condition happens because of abnormal vitreous traction on the fovea. As the vitreous gel ages and separates from the retina, it can pull on the macula, tearing a hole in it. Associated causes include high myopia, trauma, and conditions causing abnormal vitreoretinal adhesion such as proliferative vitreoretinopathy.
\section{Risk Factors}
\begin{itemize}[noitemsep]
  \item Age over 55 years
  \item Female sex
  \item High myopia
  \item Eye trauma
  \item Prior posterior vitreous detachment or retinal surgery
  \item Cystoid macular edema
\end{itemize}
\section{Signs and Symptoms}
\subsection{Common symptoms}
\begin{itemize}[noitemsep]
  \item Blurred, distorted central vision (straight lines look wavy)
  \item A central dark or gray spot (scotoma)
  \item Difficulty reading, driving, or recognizing faces
  \item Peripheral vision is usually preserved
\end{itemize}
\subsection{Emergency warning signs}
\begin{itemize}[noitemsep]
  \item Sudden loss of central vision or a new central blind spot requires urgent eye examination
\end{itemize}
\section{Progression and Stages}
Macular holes are staged from early to full thickness.
\begin{itemize}[noitemsep]
  \item Stage 1: impending hole with cyst formation, no full-thickness defect
  \item Stage 2: small full-thickness hole develops
  \item Stage 3: larger full-thickness hole
  \item Stage 4: full-thickness hole with complete posterior vitreous detachment
\end{itemize}
Without surgery, the hole usually enlarges and central vision deteriorates progressively; up to half of stage 1 holes may resolve spontaneously with vitreous separation.
\section{Complications}
\begin{itemize}[noitemsep]
  \item Progressive central vision loss
  \item Retinal detachment (in a minority of cases)
  \item Metamorphopsia and central scotoma
  \item Reduced quality of life and independence
  \item Post-surgical complications such as elevated eye pressure or cataract
\end{itemize}
\section{Diagnosis}
\begin{itemize}[noitemsep]
  \item Dilated fundus examination
  \item Optical coherence tomography (OCT) is the definitive test, confirming the hole and its stage
  \item Visual acuity testing
  \item Watzke-Allen test (slit-lamp testing for a central scotoma) may aid diagnosis
\end{itemize}
\section{Prevention}
\begin{itemize}[noitemsep]
  \item Macular holes are largely not preventable
  \item Early recognition and prompt referral of new central visual symptoms
  \item Routine eye examinations, especially for older adults
  \item Protection from eye trauma
\end{itemize}
\section{Treatment Options}
\begin{itemize}[noitemsep]
  \item Observation with monitoring for small or early-stage holes
  \item Vitrectomy surgery with gas tamponade is the standard treatment, closing the hole in most cases
  \item Face-down positioning after surgery to hold the gas bubble against the macula
  \item Occasionally, enzymatic vitreolysis or repeated surgery for persistent holes
  \item Postoperative cataract extraction is commonly needed within a few years
\end{itemize}
\section{Living with It}
\begin{itemize}[noitemsep]
  \item Follow positioning instructions carefully after surgery to maximize hole closure
  \item Do not fly or travel to high altitude while a gas bubble is present
  \item Use magnification and good lighting while vision recovers
  \item Attend scheduled follow-up visits
\end{itemize}
\section{Outlook}
With vitrectomy surgery, the hole closes in roughly 90\% of cases, and most patients gain significant visual improvement, though full normal vision is uncommon. Early-stage holes and shorter duration of symptoms are associated with better outcomes.
